\scanpage{42}
\margin{41}
\begingroup
\newsavebox{\scanXLIIbox}
\newlength{\scanXLIIheight}
\begin{lrbox}{\scanXLIIbox}
\begin{minipage}{\linewidth}
\footnotesize
\noindent
\begin{minipage}[t]{.25\linewidth}
\vspace{0pt}
\begin{tikzcd}[ampersand replacement=\&,column sep=1.1em,row sep=3.1em,
  execute at end picture={
    \draw[line width=.25pt]
      ([xshift=3pt]\tikzcdmatrixname-2-1.north east)
      .. controls ([xshift=15pt,yshift=12pt]\tikzcdmatrixname-2-1.north east)
      and ([xshift=5pt,yshift=-18pt]\tikzcdmatrixname-1-1.south east)
      .. ([xshift=5pt,yshift=-6pt]\tikzcdmatrixname-1-1.south east);
    \coordinate (scanXLIIpatch) at ([xshift=5pt,yshift=-8pt]\tikzcdmatrixname-1-1.south east);
    \draw[line width=.55pt]
      ([xshift=-2pt,yshift=2pt]scanXLIIpatch) -- ([xshift=-1pt,yshift=4pt]scanXLIIpatch)
      -- ([xshift=1pt,yshift=4.2pt]scanXLIIpatch) -- ([xshift=-.5pt,yshift=1pt]scanXLIIpatch)
      -- ([xshift=2pt,yshift=3pt]scanXLIIpatch) -- ([xshift=3pt,yshift=1.3pt]scanXLIIpatch)
      -- ([xshift=1pt,yshift=2pt]scanXLIIpatch) -- ([xshift=2pt,yshift=-1pt]scanXLIIpatch)
      -- ([xshift=.4pt,yshift=-2pt]scanXLIIpatch) -- ([xshift=2.5pt,yshift=-2.5pt]scanXLIIpatch);
    \draw[line width=.35pt]
      ([xshift=1pt,yshift=2pt]scanXLIIpatch) .. controls ([xshift=-1pt,yshift=-1pt]scanXLIIpatch)
      and ([xshift=2pt,yshift=-2pt]scanXLIIpatch) .. ([xshift=2pt,yshift=-.5pt]scanXLIIpatch);
    \draw[line width=.3pt]
      ([xshift=3pt]\tikzcdmatrixname-2-1.north east)
      .. controls ([xshift=7pt,yshift=4pt]\tikzcdmatrixname-2-1.north east)
      and ([xshift=8pt,yshift=12pt]\tikzcdmatrixname-2-1.north east)
      .. ([xshift=6pt,yshift=11pt]\tikzcdmatrixname-2-1.north east)
      .. controls ([xshift=3pt,yshift=8pt]\tikzcdmatrixname-2-1.north east)
      and ([xshift=6pt,yshift=4pt]\tikzcdmatrixname-2-1.north east)
      .. ([xshift=3pt]\tikzcdmatrixname-2-1.north east);
  }]
f^{-1}(U)\arrow[rr,hook,"i"]\arrow[d,"f"]\&\&X\arrow[d,"f"]\\
U\arrow[urr,"g"']\&\underset{j}{\unsure{\rightarrow|\mathbin{-}}}\&Y
\end{tikzcd}
\par\medskip
\[
\begin{gathered}
gf=i\\[.7em]
fg=j
\end{gathered}
\]
\end{minipage}\hfill
\begin{minipage}[t]{.72\linewidth}
\vspace{2.3ex}
$f$ morphisme birationnel
\makebox[0pt][l]{\raisebox{2.0ex}{\inserted{propre}}}%
\begin{tikzpicture}[overlay,baseline=0pt]
\draw[line width=.3pt] (.1em,.4ex) .. controls (.2em,1.8ex) and (.3em,2ex) .. (.9em,2.1ex);
\end{tikzpicture}%
de \unsure{préschémas} intègres $X,Y$\\[1.0ex]
$U$ ouvert de $Y$ sur lequel $f^{-1}$ est défini\\[.7ex]
\hspace*{.12\linewidth}[Si $Y$ est régulier en codim 1, \unsure{alors}\\[.7ex]
\hspace*{.16\linewidth}$U$ est le \unsure{complémentaire} d'un \unsure{ens}\\[.7ex]
\hspace*{.19\linewidth}de codim $\geqslant2$], \deleted{i.e. \unsure{ici|aussi} le \unsure{ft|$f$}}\\[.7ex]
\hspace*{.04\linewidth}i.e. on sait \unsure{$f^{-}$|qu'} \unsure{expr|existe} un voisinage $V$\\[.7ex]
\hspace*{.04\linewidth}tel que $f$ \unsure{induise} un \unsure{isom} $f^{-1}(U)\unsure{\rightarrow|\mathbin{-}}V$.
\end{minipage}
\par\bigskip\medskip
\noindent\hspace*{.12\linewidth}%
\begin{minipage}{.88\linewidth}
Soit $\xi\in\operatorname{Pic}(X)$ défini par un\hspace{.4em}%
\makebox[0pt][l]{\raisebox{3.2ex}{\hspace{.6em}\inserted{\uline{cycle}}}}%
\begin{tikzpicture}[overlay,baseline=0pt]
\draw[line width=.35pt] (-.25em,3.7ex) .. controls (.1em,2.9ex) and (.1em,1.0ex) .. (.25em,-.4ex)
  .. controls (.35em,1.5ex) and (2.4em,2.4ex) .. (3.5em,3.2ex);
\end{tikzpicture}%
\hspace{.4em}\uline{divisoriel} $\unsure{D|\Delta}$ ($X$ \unsure{connexe|normal})\\[1.2ex]
considérons \deleted{\illegible{1}} \deleted{le morphisme} $g^{-1}(\unsure{\eta|\xi})\in\operatorname{Pic}(U)$, il\\[1.0ex]
correspond à un cycle divisoriel
\end{minipage}
\end{minipage}
\end{lrbox}
\setlength{\scanXLIIheight}{\dimexpr\ht\scanXLIIbox+\dp\scanXLIIbox\relax}
\noindent\usebox{\scanXLIIbox}\kern-\wd\scanXLIIbox
\raisebox{-\dp\scanXLIIbox}[0pt][0pt]{%
\begin{tikzpicture}[x=\wd\scanXLIIbox,y=\scanXLIIheight,baseline=0pt]
\path[use as bounding box] (0,0) rectangle (1,1);
\draw[line width=.35pt] (.355,1.04) -- (.035,.43);
\draw[line width=.4pt] (.705,.98) -- (.245,.005);
\draw[line width=.35pt] (.855,.655) .. controls (.885,.615) and (.85,.57) .. (.81,.52) -- (.718,.35);
\draw[line width=.2pt] (.716,.346) -- (.683,.29) (.679,.283) -- (.639,.215)
  (.635,.208) -- (.598,.145) (.594,.138) -- (.55,.062) (.546,.055) .. controls (.532,.031) and (.521,.013) .. (.520,-.004);
\end{tikzpicture}}
\par
\endgroup
